\chapter{Yoshida's Hermitian forms, Weil positivity and the Thorin programme}\label{chap:yoshida}\label{yo:chapter}\chaptermark{Yoshida's Hermitian forms and Weil positivity}

% CHECK: review-A.md Section 6: the account of Yoshida's paper (Theorem 1 as positivity for a up to (1/2) log 2, Theorem 2 as a reduction of RH to nondegeneracy on a completion, and the statements (Y1) to (Y7) of Section yo:sec:yoshida) is the first author's reading and has not been checked against the paper.
In 1992 Hiroyuki Yoshida studied Weil's explicit formula as a Hermitian form on functions supported in a window $[-a,a]$ \cite{Yoshida92}. As the first author reads the paper, Yoshida proved the form positive definite for $a\le\frac12\log2$, the range in which no prime power enters, and reduced the Riemann Hypothesis to its nondegeneracy on a completed space for every $a$ (Section~\ref{yo:sec:yoshida}). This chapter reads Yoshida against the Thorin programme of this book (Section~\ref{ch:intro}) and the sine squared lemma \eqref{eq:sin2}, and works the open problems that come out of the comparison.

Yoshida's functional $\mathcal Y(F)=\sum_\rho\Phi(\rho)$ is Weil's distribution, and the heat trace $W$ of the reciprocal Xi law is half its value on a Gaussian. The sine squared lemma is the hinge: its Brownian level $x$ is Yoshida's support length $2a$. There is an uncertainty barrier for one coefficient: Proposition~\ref{yo:prop:barrier} bounds the off line part of the tent coefficient $D(x)=x\langle\chi_0,\chi_0\rangle$ at $a=x/2$ by an explicit $\varepsilon_H(x)$; against the lower bound for the on line part that the first $1700$ zeros give (a double precision computation), this keeps $D$ positive at levels below about $47$ (Proposition~\ref{yo:prop:Dpos}) whatever the zeros above the verified height $3\times10^{12}$ \cite{PT21} do, so that no computation of $D$ at those levels can detect an off line zero.

\paragraph{What is not claimed.} Nothing here proves RH. A number is computed in floating point at the precision stated, or in ball arithmetic where the text says so. The barrier concerns the tent coefficient only. The account of Yoshida's paper in Section~\ref{yo:sec:yoshida} is the first author's reading of it, used here as context; it was not checked against the paper.
